\section{Results}\label{sec:results}

\subsection{Invariance and the scope of the rapid-review limit}\label{sec:invariance-results}

Under the admissibility conditions in Section~\ref{sec:structural-methods}, the interval \(N\geq0\), \(Z\geq0\), \(0\leq E\leq E_{\max}\) is forward invariant. The logistic boundary has zero population growth at \(N=0\); the signal derivative cannot point outward through \(Z=0\) because its forcing is nonnegative; and projection keeps effort in its prescribed bounds. This is a state-space result, not a stability result.

As \(\Tr\to0\), the held flow and contemporaneous update converge on fixed finite horizons to the corresponding continuous-time system when the vector field is regular and no nonsmooth boundary is crossed. For a fixed positive \(\Tr\), local stability is decided by the eigenvalues of the complete review map, including the plant flow and update law. It is therefore possible for the rapid-review approximation to be accurate at short horizons while a finite-interval map has a different stability classification. The formal statement and proof are in Supplementary S2.

\subsection{Response windows and operator-specific crossings}\label{sec:crossings}

The archived trajectory-classified stage model reports persistent oscillation near a 3--4-year anchovy review interval and across a 6--12-year sprat interval for its baseline parameterisation. These are provisional classifications: the original solver and initial histories were not deposited. A newly specified, preregistered stage-map reconstruction did not reproduce those anchovy or sprat oscillation criteria; it did reproduce annual-review convergence for anchovy and convergence for the cod and slow-stock criteria. This mismatch is not averaged away: the reconstruction and archived classification are different evidence objects, and a mismatch does not adjudicate the unavailable legacy generating object (Supplementary S9.2).

On the common illustrative logistic plant, the exponential frozen-assessment mobilising update has a single complex-pair unit-circle crossing at \(\Tr=6.50129495\) years. The controller update exactly integrates the effort-rate ODE with the endpoint assessment fixed; the multiplier calculation uses its local linearization. The crossing angle is 0.175494 radians, or about 0.05586\(\pi\). The sampled multipliers are outside the unit circle immediately below the crossing and inside immediately above it. The forward-Euler implementation, on the same plant and baseline vector, has a complex-pair crossing at 47.535982 years and a second, real \(-1\) crossing at 79.142688 years. Thus the transition pattern is update-operator dependent; the Euler values are not estimates of an institutional threshold. A continuous-delay comparator from the companion analysis gives crossings near 3.666 and 150.358 years (Abaee, 2026), further illustrating that a continuous lag, a frozen-assessment controller update, and a forward-Euler command step are different models (Fig.~\ref{fig:rho}; the common-plant scope is detailed in Supplementary S11).

The exponential frozen-assessment protective comparator has no crossing on the tested \([0.2,200]\)-year grid, with maximum sampled spectral radius 0.996740. This numerical statement is specific to the stated gain convention, equilibrium, projection-inactive local map, and grid. It is not a general theorem that protective policies are stable.

\begin{table}[htbp]
\centering\small
\caption{Baseline local stability record on the illustrative logistic plant. A crossing is a numerical unit-circle crossing of the specified map; it is not a field-estimated review interval.}
\label{tab:crossings}
\begin{tabularx}{\linewidth}{@{}p{0.29\linewidth}p{0.31\linewidth}X@{}}
\toprule
Controller/update & Computed transition & Interpretation \\
\midrule
Mobilising, exponential frozen-assessment & \(6.501295\) yr; one complex pair & Unstable below, stable above the crossing on the scanned branch. \\
Mobilising, forward Euler & \(47.535982\) yr and \(79.142688\) yr & Complex-pair crossing, then real \(-1\) crossing; distinct Euler map. \\
Protective, exponential frozen-assessment & No crossing on \([0.2,200]\) yr; max \(\Rho=0.996740\) & Stable on the tested grid only. \\
Continuous-delay comparator & About \(3.666\) and \(150.358\) yr & Different continuous-lag operator; shown for contrast. \\
\bottomrule
\end{tabularx}
\end{table}

\begin{figure}[htbp]
\centering
\includegraphics[width=0.98\linewidth]{fig_rho_scan_v53.png}
\caption{Baseline spectral-radius scans for exponential frozen-assessment and Euler mobilising and protective maps. The exponential mobilising crossing is near 6.50 years; the plotted vertical range is restricted to \(0.9\leq\Rho\leq1.4\) for readability. Curves are model outputs, not empirical response curves.}
\label{fig:rho}
\end{figure}

\subsection{Parameter sensitivity and joint perturbations}\label{sec:sensitivity}

The baseline exponential-update crossing is fragile to some scalar parameters. In the archived one-at-a-time \(\pm5\%\) battery, increasing \(K\), \(q\), or \(E_{\max}\) can move the crossing to about 13.8--14.3 years, while decreasing these parameters can remove the crossing from the tested interval; perturbing \(\eta\) produces crossings from about 2.19 to 8.72 years. Changes to the offset \(d_0\) move the crossing only modestly, from about 6.47 to 6.53 years, while the reference scale and memory time have smaller but non-negligible effects. The crossing-only summaries of the archived \(\pm5\%\) and \(\pm10\%\) battery are tabulated in Supplementary S11.4; the full archived rows and fields are retained in \path{src/p7_empirical/recovered/sens_table.csv}. The separate rerun \(\pm2\%\) crossing table is in Supplementary S11.3.

A separate joint-box analysis perturbs six parameters \((r,K,q,\eta,E_{\max},\Delta_{\rm ref})\) simultaneously. At \(\pm0.5\%\), 60 of 64 corners have a crossing, with the first-crossing values spanning 0.871965--10.670425 years; at \(\pm1\%\), 44 of 64 have a crossing, spanning 4.423192--13.653242 years. All 4 and 20 no-crossing corners, respectively, remain inside the unit circle at every point on the 20,001-point review-interval grid. The largest sampled radii among these no-crossing corners are 0.9999713 and 0.9999953. These are dense-grid numerical classifications, not proofs of stability between grid points, and the perturbation boxes are not confidence regions. They show that the 6.50-year baseline value should not be read as a robust threshold.

\begin{figure}[htbp]
\centering
\includegraphics[width=0.87\linewidth]{fig_joint_sensitivity_v53.png}
\caption{Six-parameter \((r,K,q,\eta,E_{\max},\Delta_{\rm ref})\) joint-box sensitivity of the exponential-update first crossing. Points show the minimum and maximum crossing among corners that cross; the annotated no-crossing corners were stable at all 20,001 tested review intervals. Boxes are deterministic perturbation sets, not probability or confidence intervals.}
\label{fig:joint-sensitivity}
\end{figure}

\subsection{Forage-fish screen: no multiplicity-adjusted target-band detections}\label{sec:screen}

Among 42 RAM Legacy v4.66 forage-fish stocks with at least 20 annual SSB points, no stock-band cell is significant after Benjamini--Hochberg adjustment across 84 tests in the baseline 2.5--5-year and 5--9-year bands. The smallest raw empirical \(p\)-value is 0.0448, but its adjusted value is 1.00. A wider sensitivity battery varying detrending, end-point trimming, red-noise simulation counts, seeds, and band definitions also yields no adjusted target-band detections; one nominal baseline band-B flag does not survive multiplicity control. The result is evidence about this particular diagnostic and cohort, not evidence that governance feedback is absent.

The screen's null-excess distribution has many values below its AR(1) 95\% thresholds, but an AR(1) null controls only the fitted red-noise structure. It does not remove quasi-periodic external forcing, regime shifts, observation error, or assessment changes. In particular, the stock screen does not estimate the sign of a harvest-control rule or the causal effect of its review interval. A separate 17-cell synthetic power table detects sprat-class signals at 1.00 in the low-noise \(H=100,200\) cells and at 0.24--0.70 in the corresponding high-noise cells; anchovy-class power is 0.00--0.10 and the cod-labelled false-positive rate 0.02--0.08. These are conditional simulation proportions, not per-stock power guarantees, and the table was assembled from archived outputs rather than a complete driver rerun (Supplementary S14.1). Fig.~\ref{fig:screen} displays the baseline adjusted values and stock-level null comparisons; full per-stock results and variants are in Supplementary S9.

\begin{figure}[htbp]
\centering
\includegraphics[width=0.98\linewidth]{fig_screen_v39.png}
\caption{Baseline SSB periodicity screen. Left: Benjamini--Hochberg q-values for 84 stock-by-band tests; the dashed line denotes 0.05. Right: per-stock band power relative to the fitted AR(1) 95\% null threshold (A: 2.5--5 years; B: 5--9 years). A high ratio is a diagnostic exceedance measure, not a causal governance signal.}
\label{fig:screen}
\end{figure}

\subsection{Review-interval proxy test: no approach gradient is identifiable}\label{sec:tr-results}

The re-executed \(\Tr\)-ranked test contains 42 stocks but only two assigned proxy values: seven Peru/Chile cases at 0.5 years and 35 annual-TAC cases at 1 year. Both proxy values lie below both sampled response windows (4 years for the anchovy-class window and 8 years for the sprat-class window). There are no documented responsive multi-year cases in this cohort from which to estimate how spectral excess changes as \(\Tr\) approaches a modeled window. The correct inference is therefore \emph{no identifiable window-approach gradient}, not a supported gradient and not a direct empirical confirmation of a particular threshold.

The two-group contrasts do reproduce: the 0.5-year Peru/Chile group has greater median band-A null-excess than the annual group (0.529 versus 0.257; two-sided Mann--Whitney \(p=0.007\)), and the rank association with assigned \(\Tr\) is negative for band A (Spearman \(\rho=-0.408, p=0.007\)) and band B (\(\rho=-0.340, p=0.028\)). The sign is opposite to a simple approach-to-window prediction: the group assigned the shorter interval is farther from either modeled window. These reproduced two-group associations must not be described as an approach gradient or as an institutional effect. The 0.5-year assignments are group-level proxies for in-season responsive closures, not primary-verified stock-specific review intervals; all seven are Peru/Chile cases. Regional ecology, including ENSO-related quasi-periodic forcing, is a plausible confound, but is not established as the cause by this analysis. The AR(1) null absorbs fitted red noise, not quasi-periodic forcing. At nominal 95\% thresholds, the test records 0/42 band-A flags and 1/42 band-B flag; these are not additional BH-adjusted discoveries. Per-stock intervals, thresholds, and coding notes are in Supplementary S10.

The companion F-placement audit locally reproduces the 42-stock cohort, matches fisheries-ADH F to last-F for 42 of 43 intersected stocks, and matches all 13 broader-pool taxonomic-group F quartiles. Mean stock-level F is associated with band-A null-excess before excluding Peru/Chile (\(\rho=0.3069,p=0.0609\)) but not after exclusion (\(\rho=0.0339,p=0.8565\)); band-B results show the same attenuation. This is consistent with regional confounding, not proof of its cause. The archived ADH formula and cohort summaries also reproduce, but ADH remains a calculated buffer-years proxy rather than a direct institutional-response measure (Supplementary S13.2).

\subsection{Northern cod: a descriptive two-window record}\label{sec:cod-results}

The assessment record separates a 1991--1995 decline from the post-1995 low-stock period. SSB falls from about 730 kt in 1991 to 11 kt in 1995; the later assessment values remain low through the mid-2000s before increasing to about 300 kt in 2015. Natural mortality varies substantially across assessment years and is not a monotone proxy for governance. The window split is descriptive: it does not identify a switch in institutional timing, nor adjudicate environmental, predation, model, and management explanations. The constrained-mortality formulation supplies useful hypotheses but its underlying quantities are not reproduced in this edition. Fig.~\ref{fig:cod} and Supplementary S13.3 preserve the values and period definitions.

\begin{figure}[htbp]
\centering
\includegraphics[width=0.91\linewidth]{fig_cod_v39.png}
\caption{Northern cod assessment SSB and reported natural mortality by assessment year, with the 1991--1995 crash interval separated from the later record. The partition is descriptive and does not estimate a governance effect.}
\label{fig:cod}
\end{figure}

